\section{Reformulación del problema desde una perspectiva agent-first}

\subsection{Del centrado en el modelo al centrado en el sistema}

El curso propone la perspectiva ``agent-first'': considerar al agente como un sistema de decisión que se ejecuta de manera continua, y no como una interfaz de preguntas y respuestas de un solo uso. Este sistema contiene, como mínimo, cinco elementos básicos: representación del estado, espacio de acciones, mecanismo de memoria, mecanismo de planeación y un bucle de retroalimentación externa. Solo cuando estos cinco elementos funcionan de manera coordinada puede el agente salvar la distancia entre las tareas cortas y las tareas largas.

\begin{knowledgebox}{Abstracción mínima del agente como sistema}
El agente de lenguaje se puede simplificar como el siguiente ciclo:
\[
\text{Observe} \rightarrow \text{Reason} \rightarrow \text{Act} \rightarrow \text{Update Memory} \rightarrow \text{Replan}
\]
Cada eslabón se puede reforzar de manera independiente, pero el desempeño del sistema depende de la calidad de las interfaces entre los eslabones y del costo de latencia.
\end{knowledgebox}

\subsection{Contexto histórico y evolución de paradigmas}

Desde los primeros prototipos de AutoGPT hasta los sistemas actuales, que combinan ReAct, uso de herramientas y autorreflexión, la evolución del paradigma no ha sido ``un algoritmo que sustituye a otro'', sino más bien ``una capacidad que se suma a otra''. El curso subraya tres vías de crecimiento: el crecimiento de la calidad del razonamiento, el crecimiento de la confiabilidad de la memoria y el crecimiento de la escala de la planeación. En conjunto, determinan el límite superior de las tareas resolubles.

\begin{table}[H]
\centering
\begin{tabular}{p{2.5cm}p{4.2cm}p{4.7cm}}
\toprule
\textbf{Etapa} & \textbf{Capacidad representativa} & \textbf{Cuello de botella principal} \\
\midrule
Prototipos tempranos de agentes & Prompt + llamadas a herramientas & Propensión a la deriva y a bucles infinitos en tareas de cadena larga \\
Etapa intermedia de ingeniería & ReAct + reflexión + recuperación & Contaminación del estado, costo elevado, ajuste de parámetros complejo \\
Etapa actual de sistematización & Acciones de razonamiento + arquitectura de memoria + replaneación & La calidad del verificador y la generalización entre escenarios siguen siendo insuficientes \\
\bottomrule
\end{tabular}
\caption{La evolución de ingeniería de los agentes de lenguaje no es un avance puntual, sino una actualización simultánea de los elementos del sistema}
\end{table}

\begin{warningbox}{Malentendido frecuente: tratar al agente como un ``chatbot que sabe llamar herramientas''}
Esta clase de definición pasa por alto la gestión del estado a largo plazo y la retroalimentación de la planeación. Un sistema de agente auténtico debe manejar de manera explícita la recuperación de fallos, el cambio de estrategia, la actualización de la memoria y el restablecimiento de objetivos; de lo contrario, su capacidad se quedará limitada a tareas cortas.
\end{warningbox}

\subsection{Resumen de la sección}

La perspectiva agent-first transforma la pregunta de ``si el modelo sabe responder'' en ``si el sistema puede completar la tarea de manera estable''. Esto ofrece un criterio de evaluación unificado para los tres grandes módulos de capacidades que siguen.
